\subsection{Posterior-predictive density-ratio pricing kernels}
\label{subsec:pk-density-method}

We estimate a daily sequence of seven-calendar-day pricing kernels for Bitcoin.
For conditioning date $t$, the common state variable under both measures is the
cumulative BTC/USD log return
\begin{equation}
    R_{t,t+\tau}=\log(S_{t+\tau}/S_t), \qquad \tau=7.
\end{equation}
The horizon $\tau$ is part of the estimand: changing it changes the two
conditional return distributions, their smoothing bandwidths and common
support, and hence the density ratio. The USD-money-market numeraire is used
for the risk-neutral distribution. With the maintained zero interest and
dividend rates, the pricing kernel is therefore estimated as
\begin{equation}
    \widehat m_t(r)=
    \frac{\widehat q^{\mathrm{USD}}_t(r)}{\widehat p_t(r)}.
    \label{eq:pk-density-ratio}
\end{equation}

Our baseline implementation, denoted Scheme B, propagates physical-measure
parameter and state uncertainty but uses a daily plug-in risk-neutral
calibration. Under $P$, we combine the four MCMC chains for date $t$ only after
burn-in and thinning. Each retained SVCJ parameter vector remains paired with
the terminal latent variance from the same chain and iteration. For every
simulated path we sample one such paired posterior draw and simulate the SVCJ
law forward for seven days. Under $Q^{\mathrm{USD}}$, we condition on the
date-$t$ calibrated SVCJ parameter vector, including its initial variance, and
simulate the same seven-day cumulative BTC/USD log return. The risk-neutral
jump drift includes the compensator required by the USD martingale restriction;
we report the Monte Carlo diagnostic for
$E_{Q^{\mathrm{USD}}}[\exp(R_{t,t+\tau})]=1$.

For each measure $M\in\{P,Q^{\mathrm{USD}}\}$ and date we generate
$N=100{,}000$ paths and estimate a Gaussian kernel density,
\begin{equation}
    \widehat f_{M,t}(r;h_{M,t})=
    \frac{1}{N h_{M,t}}\sum_{i=1}^{N}
    \phi\!\left(\frac{r-R^{M}_{i,t,t+\tau}}{h_{M,t}}\right).
    \label{eq:pk-kde}
\end{equation}
MATLAB's normal-reference selector determines separate baseline bandwidths for
$P$ and $Q^{\mathrm{USD}}$. We assess smoothing sensitivity by multiplying
both measure-specific baseline bandwidths by $0.75$, $1.00$, and $1.25$. All
six densities are evaluated on the same 2,001-point grid and numerically
renormalized on that grid before division.

Density ratios are reported only on empirically supported states. Let
$I^P_t$ and $I^Q_t$ denote the respective sample intervals bounded by the
$0.5$ and $99.5$ percentiles. The valid support is
$A_t=I^P_t\cap I^Q_t$. We set $\widehat m_t(r)$ to missing outside $A_t$ and
do not extrapolate, winsorize, or clip the ratio. A date fails descriptively if
the common interval is empty or too narrow, a KDE is nonfinite, or the
physical density is nonpositive on valid grid cells. For every bandwidth we
audit
\begin{equation}
 \int \widehat p_t(r)\,dr=1,\qquad
 \int \widehat q^{\mathrm{USD}}_t(r)\,dr=1,\qquad
 \int_{A_t}\widehat p_t(r)\widehat m_t(r)\,dr
 =\int_{A_t}\widehat q^{\mathrm{USD}}_t(r)\,dr.
 \label{eq:pk-numerical-audit}
\end{equation}

The January 2024 pilot contains all 31 conditioning dates. The maximum density
normalization error is $4.44\times10^{-16}$, the maximum density-ratio identity
error is $5.55\times10^{-16}$, and the smallest baseline common-support mass is
$0.9366$. The $Q^{\mathrm{USD}}$ martingale z-statistics range from $-2.05$ to
$1.42$. These are numerical validation results, not evidence that the physical
model is final. In particular, every January physical-measure run fails the
current strict convergence screen and records 10--14 parameter warnings. We retain
and flag these dates rather than selecting on convergence while the
$P$-measure estimation is being improved.

Finally, the conditioning information sets are not exactly synchronous. The
physical rolling sample is formed from daily closing observations, whereas the
risk-neutral parameters are calibrated from Deribit option transactions near
08:00 UTC. Thus, Equation~\eqref{eq:pk-density-ratio} compares the same
BTC/USD return, horizon, and USD numeraire, but it remains subject to a clock
mismatch between the $P$ and $Q^{\mathrm{USD}}$ inputs. This limitation should
be removed or explicitly controlled before the kernels are used for formal
economic inference or live allocation.
